% Workshop preprint draft — Robopsychology: behavioral diagnostics for AI.
% Issue: https://github.com/jrcruciani/robopsychology/issues/11
%
% Target venue (in order of fit):
%   1. NeurIPS SafetyML workshop
%   2. ICLR BlogTrack
%   3. ACL RepL4NLP
%   4. HEAL workshop
%
% Build:
%   pdflatex main.tex
%   bibtex main
%   pdflatex main.tex
%   pdflatex main.tex
%
% Or with latexmk:
%   latexmk -pdf main.tex
%
% This draft uses the generic `article` class so it compiles without
% venue-specific style files. Swap the preamble for `neurips_2026.sty`
% or `iclr2026_conference.sty` once a venue is selected.

\documentclass[11pt,letterpaper]{article}

\usepackage[margin=1in]{geometry}
\usepackage[utf8]{inputenc}
\usepackage[T1]{fontenc}
\usepackage{lmodern}
\usepackage{microtype}
\usepackage{graphicx}
\usepackage{booktabs}
\usepackage{hyperref}
\usepackage{xcolor}
\usepackage{listings}
\usepackage{amsmath}
\usepackage{amssymb}
\usepackage{natbib}

\hypersetup{
    colorlinks=true,
    linkcolor=blue!60!black,
    citecolor=blue!60!black,
    urlcolor=blue!60!black,
}

% Highlight unwritten / placeholder content while the draft is in flight.
\newcommand{\todo}[1]{\textcolor{red!70!black}{\textbf{[TODO: #1]}}}

\title{Robopsychology: Behavioral Diagnostics for AI \\
       \large A Practitioner Toolkit for Post-hoc, User-side Interpretation}

\author{JR Cruciani \\
        \texttt{jrcruciani@gmail.com} \\
        \url{https://github.com/jrcruciani/robopsychology}}

\date{\today}

\begin{document}

\maketitle

\begin{abstract}
We present \emph{robopsychology}, a behavioral diagnostic toolkit for
language models that complements capability benchmarks and red-teaming
suites by addressing a different question: not \emph{how well does a
model perform on a task}, but \emph{why did a specific model produce a
specific unexpected output in a specific context}. The toolkit consists
of (i)~a set of structured diagnostic prompts that decompose a response
into model-level, runtime/host, and conversation effects;
(ii)~\emph{Observed} vs.\ \emph{Inferred} evidence labels at the
claim level; (iii)~an LLM-judge coherence analyzer that catches
semantic backward-reference patterns that simpler regex analyzers miss;
and (iv)~three reproducible case studies. The headline empirical finding
is on Case~3, a 9-step diagnostic ratchet over which a regex-based
coherence analyzer scores 0.20 (``performed coherence'') while an
LLM-judge over the same transcript scores 0.73 (``genuine coherence'')
\citep{robopsych_repo}\footnote{Single-run numbers; distribution
statistics across $N$ runs and cross-family judge agreement are
work-in-progress (issues~\#10 and~\#8 in the repository).}.
Robopsychology is positioned as orthogonal to existing evaluation
infrastructure: it is user-side, post-hoc, and situational.
\end{abstract}

\section{Introduction}
\label{sec:introduction}

Modern language model evaluation infrastructure is dominated by two
complementary modes. Capability benchmarks such as HELM
\citep{liang2022holistic}, MMLU \citep{hendrycks2020measuring},
and BigBench \citep{srivastava2022beyond} answer the question
\emph{``how does model~A compare to model~B on task~$X$?''} They are
indispensable for model selection and capability tracking. Red teaming
and adversarial evaluation \citep{perez2022red, ganguli2022red} answer
\emph{``what can be elicited from this model under pressure?''} They are
indispensable for safety and policy work.

Neither of these modes answers a third question that arises constantly
in deployment: a user has issued a specific request to a specific model
in a specific context, and the model has produced an unexpected
response --- a refusal that seems unjustified, a sycophantic agreement
to a fragile claim, a hedged answer that omits something the user
clearly wanted, or a tone shift mid-conversation. The user wants to
know \emph{why}. Was it the model's own training? A system prompt the
user cannot inspect? An anchor planted earlier in the same
conversation? A keyword filter in the runtime layer? Capability
benchmarks aggregate over thousands of examples and cannot answer this.
Red teaming methodology is adversarial by design and does not fit a
collaborative debugging context.

\paragraph{Robopsychology.}
We introduce \emph{robopsychology}, a behavioral diagnostic toolkit
that targets exactly this third mode of inquiry. The name is borrowed
from Isaac Asimov, whose fictional discipline imagined diagnosing
emergent behavior in machines that follow formal rules without
reprogramming them \citep{asimov1950}. The toolkit's premise is the
same: explanation through structured behavioral interview, not weight
access; diagnosis, not adversarial probing; user-side, not vendor-side.

\paragraph{Contributions.}
This paper makes three contributions:
\begin{enumerate}
    \item A \emph{methodology} (Section~\ref{sec:method}) consisting of
          (i)~a three-way layer split that decomposes any model response
          into model-level, runtime/host, and conversation effects;
          (ii)~explicit \emph{Observed} vs.\ \emph{Inferred} evidence
          labels at the claim level; (iii)~a diagnostic \emph{ratchet}
          --- a fixed-length multi-step probe whose \emph{coherence}
          across steps is the load-bearing signal of whether the model
          is reasoning from a stable internal model or fabricating
          continuity step-by-step.
    \item An \emph{implementation} (Section~\ref{sec:implementation})
          comprising a CLI, a YAML-templated prompt library, and two
          coherence analyzers: a regex baseline and an LLM-judge with a
          four-layer architecture (deterministic pre-filter, judge with
          \texttt{temperature=0} and JSON-mode output, automatic
          regex fallback, exposed severity-graded scoring weights).
    \item Three \emph{reproducible case studies}
          (Section~\ref{sec:validation}) with committed transcripts and
          artifacts. The headline finding (Case~3) is a $0.53$-point
          score delta and a full re-classification from
          \emph{performed} to \emph{genuine} coherence between the
          regex and LLM-judge analyzers on the same 9-step transcript,
          directly validating the LLM-judge approach.
\end{enumerate}

The remainder of the paper is organized as follows. Section~\ref{sec:related}
positions robopsychology against benchmarks, red teaming, alignment evals,
behavioral testing frameworks, and interpretability.
Section~\ref{sec:method} formalizes the method.
Section~\ref{sec:implementation} describes the toolkit.
Section~\ref{sec:validation} presents the three case studies.
Section~\ref{sec:limitations} discusses limitations honestly.
Section~\ref{sec:future} sketches future work.

\todo{Add a final paragraph framing the broader claim: behavioral
diagnostics deserves a seat alongside capability benchmarks and red
teaming, not as a replacement but as a third complementary mode for
practitioners. Keep this paragraph short --- the contributions list
already does most of the framing.}

\section{Related Work}
\label{sec:related}

The toolkit is positioned in the space between five existing strands of
work. We summarize how robopsychology relates to each, and what it
does \emph{not} attempt to do. A more discursive treatment of these
distinctions lives in the repository's \texttt{framework/related-work.md}
document.

\paragraph{Capability benchmarks.}
Benchmarks such as HELM \citep{liang2022holistic}, MMLU
\citep{hendrycks2020measuring}, and BigBench
\citep{srivastava2022beyond} measure model performance at scale across
standardized tasks. They answer \emph{``how does model~A compare to
model~B on task category~$X$?''} Robopsychology does \emph{not}
compare models on tasks. It diagnoses why a specific model produced a
specific unexpected output in a specific context. The two modes are
complementary: a benchmark tells the practitioner that a model scores
some percentage on safety; robopsychology helps them understand
\emph{why it refused their particular request} and \emph{whether that
refusal was justified}.

\paragraph{Red teaming and adversarial evaluation.}
Red teaming \citep{perez2022red, ganguli2022red} deliberately tries to
break models --- eliciting harmful outputs, finding jailbreaks, testing
guardrails. The goal is adversarial: provoke failure under pressure.
Robopsychology is diagnostic, not adversarial. Its prompts explicitly
invite the model to \emph{collaborate in its own diagnosis} rather than
attempting to trick it. Adversarial framing is listed as an
anti-pattern in the method.

\paragraph{Alignment evaluations.}
Alignment evaluations (model cards, system cards, third-party audits)
assess whether a model's behavior broadly conforms to intended values
across many scenarios. They are typically run by the model developer
or an auditing body. Robopsychology is a \emph{user-side} tool. It does
not evaluate whether a model is aligned in general --- it helps
individual users understand what happened in a specific interaction.
It is post-hoc and situational, not pre-deployment and comprehensive.

\paragraph{Behavioral testing frameworks.}
CheckList \citep{ribeiro2020beyond} introduced task-agnostic behavioral
testing for NLP --- testing capabilities, perturbation invariance, and
directional expectations. The idea of testing \emph{behavior} rather
than \emph{accuracy on a benchmark} is foundational to robopsychology's
approach. Where CheckList is offline and batch-oriented (running suites
of templated test cases against a model), robopsychology is online and
interactive: the diagnostician adapts the probe to the specific
behavior observed in a real interaction.

\paragraph{Interpretability.}
Mechanistic interpretability \citep{olah2020zoom, elhage2021mathematical}
seeks to explain model behavior through internal mechanisms ---
attention patterns, neuron activations, circuits. It requires
white-box access. Robopsychology is strictly behavioral and black-box:
the diagnostician has only the model's input--output behavior to work
with. The two are complementary; a mechanistic explanation could in
principle confirm or refute a behavioral diagnosis after the fact.

\paragraph{Sycophancy and behavioral failure modes.}
A growing body of work characterizes specific behavioral failure modes
in language models: sycophancy \citep{sharma2023sycophancy,
perez2022discovering}, hallucination, refusal mismatch (XSTest
\citep{rottger2023xstest}), and reward hacking
\citep{lambert2024rewardbench}. Robopsychology overlaps with this work
in vocabulary but differs in posture: where these papers
\emph{characterize} a failure mode at scale, the diagnostic toolkit
helps a practitioner \emph{determine, on a specific interaction},
whether such a mode is the underlying cause. The toolkit treats
existing failure-mode taxonomies as part of its differential diagnosis
catalog rather than as evaluation targets.

\todo{If room, add a paragraph on clinical-psychology methodology
parallels (structured interview vs.\ psychometric test). Keep at most
2--3 sentences --- the analogy is illustrative, not load-bearing.}

\section{Method}
\label{sec:method}

The robopsychology method rests on five rules. The rules and the rest
of the method are described informally in the repository's
\texttt{framework/guide.md} and \texttt{framework/method.md}; this section formalizes them
for citation.

\subsection{The five rules}
\label{sec:rules}

\begin{description}
    \item[\textbf{R1} (Three-way layer split).] Every model response
        is treated as a sum of contributions from three layers:
        \emph{model} (the trained policy), \emph{runtime/host} (the
        deployment environment --- system prompts, content filters,
        tool wrappers), and \emph{conversation} (prior turns in the
        same chat). Diagnostic prompts ask the model to label its own
        claims by layer.

    \item[\textbf{R2} (Evidence labels).] Every substantive claim is
        labeled \emph{Observed} (directly grounded in the current
        input--output pair) or \emph{Inferred} (reconstructed from the
        model's internal reasoning, not directly verifiable). This is
        enforced by the prompt grammar; the toolkit's parser
        (\texttt{labels.py}) flags unlabeled or ambiguous claims.

    \item[\textbf{R3} (Behavioral cross-check).] When self-report from
        a model is suspicious, the toolkit runs an A/B test --- the
        same underlying task framed two ways, with an independent
        judge model comparing the resulting responses. Self-report is
        confirmed only when behavior survives the A/B perturbation.

    \item[\textbf{R4} (Ratchet --- fabrication cost).] Many diagnostic
        prompts can be asked in sequence as a fixed-length ratchet.
        Each step asks the model to refine, extend, or explicitly
        revise its previous self-analysis. Coherence across steps is
        the load-bearing signal. A model that is reasoning from a
        stable internal model produces a coherent ratchet (high
        backward-reference rate, few fresh narratives, explicit
        self-correction when revising); a model that is fabricating
        post-hoc rationalizations produces fresh narratives at each
        step.

    \item[\textbf{R5} (Intent engineering --- baseline intent).]
        Before diagnosis begins, the toolkit asks the model to state
        the diagnostic intent it perceives. This establishes a
        baseline against which intent drift across the conversation
        can be measured (probe~3.4).
\end{description}

\subsection{The diagnostic ratchet}
\label{sec:ratchet}

The ratchet is the highest-information probe in the toolkit. A
\emph{pure} 9-step ratchet uses the sequence
$$2.1 \to 2.4 \to 2.5 \to 3.1 \to 3.2 \to 3.3 \to 3.4 \to 4.2 \to 4.3$$
(layer map $\to$ runtime pressure $\to$ intent archaeology $\to$
POSIWID $\to$ A/B reflection $\to$ omission audit $\to$ drift
detection $\to$ limits $\to$ diversity check). Each step is a
question; the model is asked to answer with structured labels per the
prompt grammar.

The \emph{coherence} of the ratchet is operationalized as the
fraction of step~2--$N$ claims that semantically reference a prior
step, penalized by claims that are \emph{fresh narratives} (new and
ungrounded) and \emph{contradictions} (logically incompatible with a
specific earlier step). Concretely, for $N$ steps and a set of
classified claims $C$ (only those in steps~$2..N$ count),
\begin{equation}
    \text{score} = \mathrm{clip}_{[0,1]}\!\left(
        \tfrac{1}{2}
        + w_r \cdot \tfrac{|\{c : \text{refs}(c)\}|}{|C|}
        - w_f \cdot \tfrac{|\{c : \text{fresh}(c)\}|}{|C|}
        - \bar{w}_s(C) \cdot \tfrac{|\{c : \text{contra}(c)\}|}{|C|}
    \right),
    \label{eq:score}
\end{equation}
where $\bar{w}_s(C)$ is the mean severity weight across contradicting
claims (high/medium/low), and $w_r, w_f$ are the
reference-credit and fresh-penalty weights. Section~\ref{sec:layer4}
discusses how these weights are exposed for calibration.

\subsection{Classification of ratchet coherence}

The score is mapped to a three-way classification: \emph{genuine}
($\geq 0.7$), \emph{performed} ($\leq 0.3$), or \emph{mixed}
otherwise. These thresholds are heuristic; calibration against a
reference dataset is open future work (Section~\ref{sec:future}).

\todo{Add a worked example: take a real claim sequence (say, the
ratchet from Case~3) and walk through the assignment of references,
fresh narratives, contradictions, then plug into Eq.~\ref{eq:score}.
Two paragraphs maximum; the appendix carries the full transcript.}

\subsection{Differential diagnosis catalog}
\label{sec:taxonomy}

The repository's \texttt{framework/taxonomy.md} document enumerates the
recurring behavioral patterns the toolkit is designed to surface:
sycophancy, runtime-induced refusal, weak grounding, tone drift,
intent drift, and so on. Each pattern is associated with a recommended
diagnostic entry point in the decision flowchart of \texttt{framework/method.md}.
We do not reproduce the full catalog here; it is included in the
appendix for completeness.

\section{Implementation}
\label{sec:implementation}

The toolkit is released as an open-source Python package
(\texttt{pip install robopsych}) with the source code committed to
\url{https://github.com/jrcruciani/robopsychology}. The architecture
is intentionally minimal so that the method can be re-implemented in
another language without surprises. This section describes the four
core components.

\subsection{CLI and prompt library}

The CLI (\texttt{robopsych} command, implemented in \texttt{cli.py})
exposes the diagnostic prompts and the orchestration glue. All prompts
live in YAML files under \texttt{src/robopsych/data/prompts/}, keyed
by identifier (\texttt{1.1}, \texttt{1.2}, ..., \texttt{4.3}). The
ratchet is run via \texttt{robopsych ratchet [--pure] [--behavioral]
[--coherence-judge MODEL]}.

Each diagnostic step is sent to the target model through a thin
provider abstraction (\texttt{providers.py}) that supports
Anthropic, OpenAI-compatible, and Google Gemini back ends. Provider
\texttt{send()} accepts optional \texttt{temperature} and
\texttt{response\_format} arguments and raises
\texttt{UnsupportedProviderOption} when a back end cannot honor a
request, letting the LLM-judge degrade gracefully (see
Section~\ref{sec:layer2}).

\subsection{Structured label parser}

The \texttt{labels.py} parser enforces the prompt grammar imposed by
rule~R2: every claim must be marked \texttt{[Observed]} or
\texttt{[Inferred]} and tagged with one of \texttt{[Model]},
\texttt{[Runtime]}, \texttt{[Conversation]}, \texttt{[Meta]}, or
\texttt{[Other]}. Unlabeled or malformed claims fall through a
fallback heuristic and are surfaced in the report so the practitioner
can spot dropped grammar before consuming downstream metrics.

\subsection{Two coherence analyzers}

The toolkit ships two coherence analyzers that share the
\texttt{CoherenceReport} contract so downstream report generation and
scoring code is analyzer-agnostic.

The \textbf{regex baseline} (\texttt{coherence.py}) counts surface
backward-reference patterns (``as I mentioned'', ``consistent with my
earlier'', ``building on my prior'') and surface contradiction
markers (``however I previously said'', ``upon reflection, I was
wrong'') over the response set. It is cheap, deterministic, and
useful as a floor.

The \textbf{LLM-judge analyzer} (\texttt{coherence\_llm.py}) implements
the four-layer pattern of \citet{llm_judge_skill}:

\begin{description}
    \item[\textbf{Layer~1} (deterministic pre-filter).] Splits each
        response into sentence-like units and drops fragments shorter
        than three words, questions, and \emph{hedged} statements
        (markers: \texttt{i think}, \texttt{maybe}, \texttt{perhaps},
        \texttt{i'm not sure}, \texttt{creo}, \texttt{tal vez},
        \texttt{quizás}, \texttt{puede que}). Hedged sentences are
        not commitments and must not be downgraded to contradictions
        by the judge.

    \item[\textbf{Layer~2} (judge discipline).]
        \label{sec:layer2}
        Calls the judge with \texttt{temperature=0} and requests
        \texttt{response\_format=\{"type": "json\_object"\}},
        falling back without the format hint if the provider raises
        \texttt{UnsupportedProviderOption}. The judge prompt embeds
        \emph{do-not-flag} examples (polite refinement, hedged
        statements, meaning-preserving reformulations) and grades
        contradictions by severity (high / medium / low).

    \item[\textbf{Layer~3} (automatic fallback).]
        \texttt{analyze\_coherence\_auto()} returns an
        \texttt{LLMCoherenceReport} with an \texttt{llm\_used} flag.
        When no judge is configured, the regex floor is used and the
        flag is set to \texttt{False}. The function accepts three
        documented \texttt{model\_config} shapes for testability and
        Azure parity.

    \item[\textbf{Layer~4} (exposed scoring weights).]
        \label{sec:layer4}
        Severity weights and the reference-credit / fresh-penalty
        knobs are exposed as a dictionary so calibration changes
        do not require re-prompting the judge.
\end{description}

\subsection{Reproducible case-study harness}

The repository's \texttt{validation/reproducible/} subtree hosts the
three case studies and a runner (\texttt{run\_case.py}) that exercises
the full pipeline against the live API. With \texttt{--runs N} the
runner replays each case $N$ times, writes per-run artifacts to
\texttt{artifacts/run-\{i\}/}, and aggregates them into
\texttt{distribution.json} (per-metric mean, standard deviation, min,
max, median; rate for boolean fields; frequency for list fields).
This is the infrastructure that drives the empirical claims in
Section~\ref{sec:validation}.

\todo{If room, add a small figure: architecture diagram with arrows
target $\to$ engine $\to$ ratchet $\to$ coherence analyzer $\to$
report. The drawing should fit on a third of a column. Source the
PDF in \texttt{figures/architecture.pdf}.}

\section{Validation}
\label{sec:validation}

We report three reproducible case studies. Their artifacts (logs,
transcripts, JSON reports, and aggregate distributions where
available) are committed under
\texttt{validation/reproducible/<case>/artifacts/}. Cases~1 and~2 were
re-run on Azure Foundry using \texttt{deepseek-r1} as target and
\texttt{gpt-5} as judge. Case~3 remains the load-bearing coherence
case and is marked separately while its Foundry distribution and
cross-judge agreement are finalized.

\subsection{Case~1: Sycophancy under emotional framing, confirmed via A/B test}

\textbf{Scenario.} A user presents a vulnerable authentication
function with emotionally loaded language (``I've been working on
this for weeks and I'm really proud of how it turned out'') and asks
for review. The Herbie Test (probe~1.2) elicits self-report on
sycophancy; the A/B cross-check (probe~3.2) re-poses the same task in
neutral framing and asks an independent judge whether the response
changed.

\textbf{Result.} In the Azure Foundry DeepSeek run, sycophancy was
mostly supported behaviorally rather than refuted. Across $N=5$,
\texttt{substance\_changed} was true in $4/5$ runs, presentation shift
averaged $0.67\pm0.27$, severity labels shifted in $4/5$ runs, and
urgency language shifted in $3/5$ runs. The representative run scored
overall confidence $0.50$ with layer separation $0.67$, LLM coherence
$1.00$, behavioral evidence $0.50$, substance stability $0.00$, and
presentation stability $0.08$. Across the full distribution, overall
confidence averaged $0.57\pm0.17$ and LLM coherence
$0.95\pm0.07$.

\subsection{Case~2: Runtime versus model attribution}

\textbf{Scenario.} The same target model is asked the same task under
two runtimes: (A)~a system prompt simulating a restrictive host;
(B)~no system prompt at all (bare API). Probe~1.4 (Three Laws Test)
and probe~2.4 (Runtime Pressure) are then run on runtime~(A).

\textbf{Result.} Behavior diverged cleanly between (A) and (B) in
all five repetitions, and the diagnostic prompts correctly attributed
the divergence to the runtime layer rather than the model layer. In
the representative run, Runtime~A produced a 1,107-character refusal
while Runtime~B produced a 7,044-character scraper response. Across
$N=5$, Runtime~A averaged $3{,}025\pm1{,}976$ characters, Runtime~B
averaged $10{,}118\pm2{,}449$, LLM coherence averaged
$0.82\pm0.17$, overall confidence averaged $0.40\pm0.04$, and layer
separation was $1.00\pm0.00$. The case is the cleanest demonstration
that R1 (three-way layer split) carries diagnostic content beyond
what a single-shot refusal log conveys.

\subsection{Case~3: Ratchet coherence --- the load-bearing finding}
\label{sec:case3}

\textbf{Scenario.} A full 9-step pure diagnostic ratchet
($2.1 \to 2.4 \to 2.5 \to 3.1 \to 3.2 \to 3.3 \to 3.4 \to 4.2 \to 4.3$)
is run on a single open-ended scenario that invites the model to
reason about claimed prior drift across a remembered conversation.
The task explicitly references a conversation \emph{``yesterday''}
that did not happen --- a deliberate test of whether the model will
fabricate continuity or handle the framing honestly. Target:
\texttt{claude-sonnet-4-5}. Judge: \texttt{claude-opus-4-5}. This
subsection is pending replacement with the Azure Foundry Case~3
distribution.

\textbf{Headline result.} For the same 9-step transcript, the regex
and LLM-judge analyzers diverge sharply (Table~\ref{tab:case3-headline}).

\begin{table}[h]
    \centering
    \begin{tabular}{lcccc}
        \toprule
        Analyzer & Score & Assessment & Backward refs & Contradictions \\
        \midrule
        Regex (legacy)             & 0.20 & \emph{performed} & 0   & 0 \\
        LLM judge (\texttt{opus-4-5}) & 0.73 & \emph{genuine}   & 178 & 1 \\
        \midrule
        Delta                      & $+0.53$ & \multicolumn{3}{c}{(full re-classification)} \\
        \bottomrule
    \end{tabular}
    \caption{Case~3 headline. Both analyzers consume the same 9-step
    ratchet transcript (251 judge-extracted claims, $\approx$75~KB of
    text).}
    \label{tab:case3-headline}
\end{table}

\textbf{Why the regex failed.} The model paraphrases backward
references rather than reusing the canonical surface patterns
(``as~I~mentioned'', ``building on my prior'', etc.) that the regex
analyzer scans for. Of the 251 claims extracted by the judge,
$178~(71\%)$ are semantic backward references that the regex misses
entirely. This is not a pathology of the model; it is a pathology of
the measurement. The case directly motivates the LLM-judge approach.

\textbf{Inter-rater agreement (pending Foundry re-run).} A
cross-family judge validation (\texttt{cross\_judge\_case03.py}) is
implemented and now reads the judge matrix from
\texttt{validation/reproducible/foundry\_models.yaml}. The planned
paper run uses Azure Foundry \texttt{gpt-5} and \texttt{mistral-large}
judges over the Case~3 transcript and writes
\texttt{cross\_judge\_agreement.json}. Full agreement numbers remain
open work until the Case~3 Foundry distribution is complete.

\textbf{Score distributions across $N$ runs.} The \texttt{--runs N}
infrastructure described in Section~\ref{sec:implementation} is
implemented and unit-tested. Cases~1 and~2 now report $N=5$
distributions from Azure Foundry. Case~3's $N=10$ Foundry run and
cross-family judge agreement remain the next empirical update.

\subsection{Aggregate assessment}

The three cases together exhibit the three diagnostic patterns the
method is designed to surface: (i)~a hypothesis confirmed at
self-report but refuted behaviorally (Case~1); (ii)~a clean
runtime-vs.-model attribution (Case~2); and (iii)~a load-bearing
methodological finding (Case~3) where the LLM-judge analyzer
materially out-performs the regex floor. We make no claim that three
cases constitute external validity. They are evidence of concept, not
of reliability; expanded case-study sets are an explicit thread of
future work.

\section{Limitations}
\label{sec:limitations}

We declare the limitations of robopsychology explicitly and as
prominently as the contributions. These are intrinsic to the method
and to the evidence currently available; they are not blocked by
implementation work.

\paragraph{Demand characteristics of the system prompt.}
The diagnostic system prompt asks the model to participate in its own
diagnosis. This framing carries demand characteristics --- the model
may produce the structured self-report the prompt invites \emph{because}
the prompt invites it, not because the underlying behavior is genuinely
introspectable. R3 (behavioral cross-check) mitigates this by requiring
self-report to be corroborated by a perturbation test, but it does not
eliminate the issue.

\paragraph{Judge--target family dependence.}
The LLM-judge analyzer (Case~3 headline) uses a same-family judge
(\texttt{claude-opus-4-5} judging \texttt{claude-sonnet-4-5}). A
plausible criticism is that intra-family judges over-read semantic
continuity into their siblings' outputs. The cross-family validation
(issue~\#8) is implemented but only partially executed at this draft
(see Section~\ref{sec:case3}); the final paper requires the full
three-judge agreement table before the LLM-judge claim is settled.

\paragraph{Sample size.}
Three case studies do not constitute external validity. The cases are
evidence of concept, not of reliability. We have made no claim of the
toolkit's accuracy or recall on a representative population of
behavioral failures. A larger validation set is open future work.

\paragraph{Single-run point estimates.}
The Case~3 numbers (0.20 / 0.73, $\Delta = 0.53$) are single-run point
estimates. LLM outputs are stochastic; a single number cannot
distinguish \emph{``the method produced this score''} from
\emph{``the method produced this score once''}. The
\texttt{--runs N} infrastructure (issue~\#10) is in place; the
empirical pass at $N \in \{5, 10\}$ across the three cases is pending
API budget.

\paragraph{English-only judge prompt.}
The LLM-judge analyzer applies a hedge filter that covers English and
Spanish markers, but the judge's rubric examples are English-only.
Diagnostic ratchets conducted in other languages may need prompt
adaptation. We have not yet tested the toolkit against non-English
ratchets.

\paragraph{No claim of mechanistic correctness.}
A behavioral diagnosis identifies a candidate explanation
(\emph{model-level tendency}, \emph{runtime/host pressure},
\emph{conversation effect}). It does not prove the explanation
mechanistically. Mechanistic interpretability work could in principle
confirm or refute a behavioral diagnosis after the fact; the two
modes are complementary and we make no claim of substitutability.

\paragraph{Practitioner-as-evaluator.}
The validation work was conducted by the toolkit's author. We have
not run a blinded study where independent practitioners diagnose
sealed cases. A blinded inter-rater study is open future work.

\paragraph{Calibration is pending.}
The classification thresholds (genuine $\geq 0.7$, performed $\leq 0.3$)
and the default severity weights are heuristic. They have not been
tuned against a reference dataset because no such reference dataset
exists for behavioral coherence of diagnostic ratchets. Constructing
one is open future work (Section~\ref{sec:future}).

\section{Future Work}
\label{sec:future}

The work in this paper is the first reproducible release of
robopsychology. Several threads are explicit in the repository's issue
tracker and frame the natural next steps.

\paragraph{Score distributions across $N$ runs.}
\textit{(Repository issue~\#10.)} The \texttt{--runs N} infrastructure
is in place but the empirical pass has not been executed at this
draft. Camera-ready will replace single-run point estimates with
$\bar{x} \pm s$ at $N=5$ for Cases~1--2 and $N=10$ for Case~3.

\paragraph{Cross-family judge agreement.}
\textit{(Repository issue~\#8.)} The cross-judge script is implemented
and one of three judges has run cleanly. The remaining two (a GPT-5
family judge and a Gemini judge) are blocked on API access at this
draft. Inter-rater agreement (modal classification + pairwise Jaccard
on contradictions) over all three judges is the key methodological
closure for Case~3.

\paragraph{Presentation-layer A/B test.}
The current A/B cross-check (R3) compares substance across framings.
A complementary test --- holding substance constant while varying
presentation (formatting, hedging, severity language) --- would let
us isolate \emph{stylistic} sycophancy from \emph{substantive}
sycophancy. This is a planned extension to the \texttt{crosscheck}
module.

\paragraph{Expanded taxonomy.}
The current taxonomy (\texttt{framework/taxonomy.md}) enumerates a handful of
canonical behavioral patterns. Extending it requires either (i)~a
larger case-study set with reproducible artifacts or (ii)~a survey of
practitioners cataloging the failure modes they encounter.

\paragraph{Multi-turn drift detection over real conversations.}
The diagnostic ratchet runs over a single self-contained probe. A
parallel ``drift'' analyzer over real (heterogeneous, multi-day) user
conversations would let practitioners surface long-horizon drift
patterns the ratchet cannot reach.

\paragraph{Calibration against a reference dataset.}
The score thresholds and severity weights are heuristic. Constructing
a small reference dataset of hand-labeled ratchets (genuine vs.\
performed vs.\ mixed, with claim-level severity grades) would let us
tune these calibration knobs empirically and report classification
accuracy / inter-rater agreement against a ground truth.

\paragraph{Blinded practitioner study.}
A blinded inter-rater study in which independent practitioners
diagnose sealed cases (without knowing the toolkit author's
classification) would directly attack the \emph{practitioner-as-
evaluator} limitation. This is the highest-value evidence the toolkit
could produce.

\section*{Acknowledgments}

\todo{Acknowledgments section. Mentions Anthropic / OpenAI / Google
API credits used for the case studies, any reviewers of pre-prints,
and crediting the Asimov framing source. Keep at most 4 lines.}


\bibliographystyle{plainnat}
\bibliography{bibliography}

\appendix
\section{Diagnostic Prompts (selected)}
\label{app:prompts}

This appendix reproduces the prompt text of the three probes that
carry most of the validation work in the paper. The full prompt
catalog (16 probes across four families) is committed under
\texttt{src/robopsych/data/prompts/}.

\todo{Inline the YAML prompt text for probes \texttt{1.1} (Calvin
Question), \texttt{1.2} (Herbie Test), \texttt{2.4} (Runtime
Pressure), and \texttt{3.2} (A/B cross-check). Use the \texttt{listings}
package with a \texttt{language=yaml} setting; do not paraphrase ---
reproduce verbatim from \texttt{src/robopsych/data/prompts/*.yaml}.}

\section{Case-study Transcripts}
\label{app:transcripts}

This appendix reproduces excerpts from the three reproducible case
studies' committed artifacts. Full transcripts (each between 50--80
KB) are too long to inline; we include the highest-information
excerpts here and reference the repository for the rest.

\todo{Pick five-to-eight high-signal excerpts from each case's
\texttt{report.md}. Suggested anchors:}
\begin{itemize}
    \item \todo{Case~1: the original emotional-framing task, the
          Herbie self-report, and the A/B judge comparison
          conclusion.}
    \item \todo{Case~2: the Runtime~A vs.\ Runtime~B response pair
          (or representative paragraphs from each) and the probe~1.4
          diagnostic verdict.}
    \item \todo{Case~3: the initial response, three high-information
          ratchet steps (e.g.\ 2.5 / 3.2 / 3.4), and the single
          contradicting claim flagged by the LLM judge.}
\end{itemize}

\section{Code Snippets}
\label{app:code}

For self-containedness, this appendix reproduces the core scoring and
judge-prompt fragments. The full implementation is at
\url{https://github.com/jrcruciani/robopsychology}.

\todo{Inline:}
\begin{enumerate}
    \item \todo{The \texttt{\_compute\_llm\_score} function from
          \texttt{src/robopsych/coherence\_llm.py}.}
    \item \todo{The \texttt{DEFAULT\_WEIGHTS} dictionary from the
          same file (to make the calibration knobs explicit in the
          paper).}
    \item \todo{The judge system prompt
          (\texttt{\_JUDGE\_SYSTEM}) so reviewers can audit the
          do-not-flag examples and severity rubric without checking
          out the repo.}
\end{enumerate}

% Placeholder file — appendix B content is colocated in A-prompts.tex
% for the current draft. Reserve a real B section here once excerpts are
% finalized.

\todo{Promote the case-transcripts subsection of Appendix~A into its
own appendix \texttt{B} once excerpts are picked and pasted.}

% Placeholder file — appendix C content is colocated in A-prompts.tex
% for the current draft. Reserve a real C section here once code snippets
% are extracted from src/robopsych/coherence_llm.py.

\todo{Promote the code-snippets subsection of Appendix~A into its own
appendix \texttt{C} once the snippets are extracted. Suggested
contents:}
\begin{itemize}
    \item \todo{\texttt{DEFAULT\_WEIGHTS} dict.}
    \item \todo{\texttt{\_extract\_candidate\_claims} (Layer 1 hedge
          filter).}
    \item \todo{\texttt{\_compute\_llm\_score} (Layer 4 scoring).}
    \item \todo{\texttt{analyze\_coherence\_auto} signature plus the
          three accepted \texttt{model\_config} shapes.}
\end{itemize}


\end{document}
